\chapter{Hilbert symbol}

\section{Hilbert symbol over a field}

\begin{definition}
	\label{def:Hilbert_symbol}
	\lean{hilbertSym}
	\leanok
	Let $k$ be a field and $a, b \in k$. If both $a,b$ are nonzero, the \emph{Hilbert symbol} $(a, b)_k$ is defined as
	follows: if there exist $x, y, z \in k$, not all zero, such that $z^2 = ax^2 + by^2$,
	then $(a, b)_k = 1$; otherwise, $(a, b)_k = -1$. If one of $a,b$ is zero, we set $(a, b)_k = 0$.
\end{definition}


\subsection{Basic properties of the Hilbert symbol}

The Hilbert symbol has the following properties, useful for computing it in practice.

\begin{lemma}
	\label{lem:well_defined_up_to_squares}
	\lean{hilbertSym.mul_square_eq}
	\uses{def:Hilbert_symbol}
	\leanok
	Let $a, b, a', b' \in k^\times$. Then $(a a'^2, b b'^2)_k = (a,b)_k$.
\end{lemma}

\begin{proof}
	Let $x, y, z \in k$ be such that $z^2 = a a'^2 x^2 + b b'^2 y^2$. Then, $xa', yb', z$ are not all
	zero and satisfy $z^2 = ax^2 + by^2$. Conversely, if $x, y, z \in k$ are such that
	$z^2 = ax^2 + by^2$, then $x/a', y/b', z$ are not all zero and satisfy
	$z^2 = a a'^2 x^2 + b b'^2 y^2$.
\end{proof}

The following lemma uses the definition \code{QuadraticAlgebra} from \mathlib.

\begin{lemma}
	\label{lem:eq_one_iff}
	\lean{hilbertSym.eq_one_iff}
	\uses{def:Hilbert_symbol}
	\leanok
	Let $a, b \in k^\times$. Then $(a, b)_k = 1$ if and only if there exists $t \in k(\sqrt{b})$ such
	that $a = N_{k(\sqrt{b})/k}(t)$.
\end{lemma}

\begin{proof}
	\begin{itemize}
		\item Case $b$ is a square in $k$. Then $k(\sqrt{b}) = k$, and the norm map is the identity, so
		$a$ is always the norm of itself. Let $c \in k^\times$ be such that $c^2 = b$.
		Then $c^2 = a \cdot 0^2 + b \cdot 1^2$, so $(a,b)_k=1$.
		\item Case $b$ is not a square in $k$. Then $k(\sqrt{b})$ is a quadratic extension of $k$, and
		the norm map is given by $N_{k(\sqrt{b})/k}(p + q \sqrt{b}) = p^2 - b q^2$.
		Suppose $a$ is a norm, so it satisfies $a = p^2 - b q^2$ for some $p, q \in k$.
		Then $1, q, p$ are not all zero and satisfy $ p^2 = a \cdot 1^2 + b \cdot q^2$,
		so $(a, b)_k = 1$. Conversely, if $x, y, z$ are not all zero and satisfy $z^2 = ax^2 + by^2$,
		then $x=0$, since otherwise $b$ would be a square in $k$, which contradicts our assumption.
		Then $x \ne 0$, and we have $a = (z/x)^2 - b (y/x)^2$, so $a$ is a norm.
	\end{itemize}
\end{proof}

\begin{proposition}
	\label{prop:hilbert_symbol_properties}
	\lean{hilbertSym.comm, hilbertSym.right_square_eq_one, hilbertSym.right_neg_self_eq_one,
	hilbertSym.right_one_minus_self_eq_one, hilbertSym.right_mul_eq_of_eq_one , hilbertSym.right_neg_mul,
	hilbertSym.right_minus_self_mul}
	\uses{def:Hilbert_symbol}
	\leanok
	The Hilbert symbol statisfies the following formulas:
	\begin{enumerate}[(i)]
		\item For $a, b \in k$, we have $(a, b)_k = (b, a)_k$.
		\item Let $a, c \in k^\times$. Then $(a, c^2)_k = 1$.
		\item For $a \in k^\times$, we have $(a, -a)_k = 1$.
		\item Let $a \in k$ with $a \ne 0, 1$. Then $(a, 1 - a)_k = 1$.
		\item Let $a, a', b \in k^\times$ and assume that $(a, b)_k = 1$. Then $(a a', b)_k = (a', b)_k$.
		\item Let $a, b \in k$. Then $(a, -ab)_k = (a, b)_k$.
		\item Let $a, b \in k^\times$ with $a \ne 1$. Then $(a, (1-a)b)_k = (a, b)_k$.
	\end{enumerate}
\end{proposition}

\begin{proof}
	\uses{def:Hilbert_symbol, lem:eq_one_iff}
	
	\begin{itemize}
		\item Trivial if $a$ or $b$ are zero. Otherwise, if $x,y,z$ is a solution to $z^2 = ax^2 + by^2$, then $y,x,z$ is a solution to
	$z^2 = bx^2 + ay^2$.
		\item In this case, $c^2 = a \cdot 0^2 + b \cdot 1^2$, so $(a,b)_k=1$.
		N.B. Alternatively: the same argument as in the first case of the proof of Lemma~\ref{lem:eq_one_iff} applies.
		\item In this case, $0 = a \cdot 1^2 + (-a) \cdot 1^2$.
		\item In this case, $1 = a \cdot 1^2 + (1 - a) \cdot 1^2$.
		\item We know, by Lemma~\ref{lem:eq_one_iff}, that there exists $t \in k(\sqrt{b})$ such that
		$a = N_{k(\sqrt{b})/k}(t)$. If $(a',b)_k=1$, then $a'$ is the norm of some $t' \in k(\sqrt{b})$,
		so $aa' = N_{k(\sqrt{b})/k}(t t')$, so $(aa', b)_k = 1$. Conversely, if $(aa', b)_k = 1$, then
		$aa' = N_{k(\sqrt{b})/k}(t'')$ for some $t'' \in k(\sqrt{b})$, so
		$a' = N_{k(\sqrt{b})/k}(t''/t)$, hence $(a', b)_k = 1$.
		\item Trivial if $a$ or $b$ are zero. Otherwise it follows from the previous steps, since $(a, -a)_k = 1$.
		\item Follows from the previous steps, since $(a, 1 - a)_k = 1$.
	\end{itemize}
\end{proof}
\subsection{Local Properties}
We now specialize to the case where $k$ is either $\R$ or $\Q_p$ for some prime $p$. In this case, the Hilbert symbol can be computed explicitly.

\begin{theorem}
	\label{thm:real}
	\lean{hilbertSym.real_eq}
	\uses{def:Hilbert_symbol}
	\leanok
	Let $a, b \in \R^\times$. Then $(a, b)_{\R} = 1$ if and only if $a > 0$ or $b > 0$; otherwise $(a, b)_{\R} = -1$.
\end{theorem}

\begin{proof}
	If $a>0$, then $(1,0,\sqrt{a})$ works. Similarly, if $b>0$, then $(0,1,\sqrt{b})$ works.
	If both $a$ and $b$ are negative, then $ax^2 + by^2 \le 0$ for all $x,y$, so there is no
	nontrivial solution to $z^2 = ax^2 + by^2$.
\end{proof}

The following uses the \mathlib API for $p$-adic numbers and integers, as well as the results from the files \code{HassePrinciple/Padics/Lemmas.lean} and \code{HassePrinciple/Padics/Legendre.lean}.


For the statement of the theorem in the $p$-adic case, we need the following definitions.


\begin{definition}
	\label{def:epsilon}
	\lean{PadicInt.epsilon}
	\leanok
	Let $u \in (\Z_2)^\times$. We define $\epsilon(u)$ to be the class modulo $2$ of $\dfrac{u-1}{2}$.
\end{definition}

\begin{definition}
	\label{def:omega}
	\lean{PadicInt.omega}
	\leanok
	Let $u \in (\Z_2)^\times$. We define $\omega(u)$ to be the class modulo $2$ of $\dfrac{u^2-1}{8}$.
\end{definition}



For the proof, we will need the following lemma.
\begin{lemma}
	\label{lem:nontrivialzero}
	\lean{Padic.exists_nontrivial_zero}
	\leanok
Let $p$ be a prime and let $v \in (\Q_p)^\times$. Let $x, y, z \in \Q_p$ be such that
$(x, y, z) \ne (0, 0, 0)$ and $z^2 - p x^2 - v y^2 = 0$. Then there exist
$z', y' \in (\Z_p)^\times$ and $x' \in \Z_p$ such that $(z', y', x')$ is a nontrivial solution to
the same equation.
\end{lemma}
\begin{proof}
	Let $(x,y,z)$ be any nontrivial solution. For any $\lambda \ne 0$, also $(\lambda x, \lambda y,
	\lambda z)$ is a solution. Taking $\lambda = p^a$ for an appropriate $a$, we can assume that
	$x,y,z$ are all in $\Z_p$ and at least one of them is a unit. Now, if $z$ is not a unit, then
	$px^2+vy^2$ is also not a unit, thus $y$ is not a unit. Conversely if $y$ is not a unit, then
	$z^2-px^2$ is also not a unit, thus $z$ is not a unit. This shows that if eithe $y$ or $z$ is not
	a unit, then they are both non units, but then $p^2$ would divide $px^2$, hence $x$ would not be a
	unit either, contradicting the fact that one of $x,y,z$ is a unit.
\end{proof}

We are now ready to compute the Hilbert symbol $(x,y)$ in the $p$-adic case.
We split the statement and the proof of the theorem for the $p$-adic into the cases where $p$ is odd
and where $p = 2$, and for each we first consider separately the cases where the valuations of $x$
and $y$ are both $0$, both $1$, or one of them is $0$ and the other is $1$.

Everywhere in the following, we assume that $x, y \in (\Q_p)^\times$, and denote by $v_x$ and $v_y$
their valuations, and by $u_x$ and $u_y$ their unit parts.

\begin{theorem}
	\label{thm:padic_odd_p}
	\lean{hilbertSym.padic_odd_eq, hilbertSym.padic_odd_case00, hilbertSym.padic_odd_case10, hilbertSym.padic_odd_case11}
	\uses{def:Hilbert_symbol, def:padicUnitPart, def:epsilon, def:legendreSym}
	\leanok
Let $p$ be an odd prime and let $x,y \in (\Q_p)^\times$. Then
\[(x,y) = (-1)^{v_x v_y\epsilon(p)}\left(\dfrac{u_x}{p}\right) \left(\dfrac{u_y}{p}\right),\]
where $\left(\dfrac{\cdot}{p}\right)$ is the Legendre symbol.
\end{theorem}

\begin{proof}
	\uses{lem:nontrivialzero, prop:hilbert_symbol_properties}
	Since the Hilbert symbol is well defined up to squares, we can assume that $v_x, v_y \in \{0, 1\}$.
	\begin{itemize}
		\item Suppose $v_x=v_y=0$. Then $x=u_x$ and $y=u_y$ are units. By the theorem of Chevalley--Warning,
		the equation $z^2 = u_x x^2 + u_y y^2$ has a nontrivial solution in $\F_p$, so it has a
		nontrivial solution in $\Q_p$, so $(x,y)=1$. Note: Chevalley--Warning is in \mathlib as
		\code{char_dvd_card_solutions}, in \code{Mathlib.FieldTheory.ChevalleyWarning}. We may need a
		multivariate Hensel's lemma for the lifting part.
		\item Suppose $v_x=1$ and $v_y=0$. Then $x=p u_x$ and $y=u_y$. Since by the previous case, we
		have $(u_x, u_y)=1$, we have $(x,y) = (p u_x, u_y) = (p, u_y)$. We just need to show that this is
		equal to $\left(\dfrac{u_y}{p}\right)$. Observe that a unit is a square in $\Q_p$ if and only if
		it is a square modulo $p$. Now, if $u_y$ is a square, then both terms are $1$, see
		\ref{prop:hilbert_symbol_properties}. Otherwise, both terms are $-1$. In fact,
		if $z^2-p^2x^2-u_yy^2$ has a nontrivial solution, we can assume that $z$ and $u_y$ are units, so
		we can reduce modulo $p$ to get $u_y = (zy^{-1})^2$, which is a contradiction.
		\item Suppose $v_x=v_y=1$. Then $x=p u_x$ and $y=p u_y$. This case follows from the properties of
		the Hilbert symbol and the previous case: $(pu_x,pu_y)=(pu_x, -p^2u_xu_y)= (pu_x, -u_xu_y)=
		\left(\dfrac{-u_xu_y}{p}\right)$. By the properties of the Legendre symbol, this is equal to
		$\left(\dfrac{-1}{p}\right)\left(\dfrac{u_x}{p}\right)\left(\dfrac{u_y}{p}\right)$, as claimed.
	\end{itemize}
\end{proof}

We proceed similarly for the case $p = 2$.

\begin{theorem}
	\label{thm:padic_2}
	\lean{hilbertSym.two_adic_eq, hilbertSym.two_adic_case00, hilbertSym.two_adic_case10, hilbertSym.two_adic_case11}
	\uses{def:Hilbert_symbol, def:padicUnitPart, def:epsilon, def:omega}
	\leanok
Let $x,y \in (\Q_2)^\times$. Then
\[(x,y) = (-1)^{\epsilon(u_x)*\epsilon(u_y) + v_x * \omega(u_y) + v_y * \omega(u_x)}.\]
\end{theorem}

\begin{proof}
	\uses{prop:hilbert_symbol_properties}
TODO
\end{proof}

\subsection{Global properties of the Hilbert symbol}

This section uses the \mathlib function \code{Nat.nth Nat.Prime}. The main goal is to prove
the product formula for the Hilbert symbol, which states that for every $a, b \in \Q^\times$, we have
\[\prod_{v \in V} (a, b)_v = 1.\]
Here, $V$ is the union of the set of prime numbers and the symbol $\infty$, and for each prime $p$,
the symbol $(a, b)_p$ is the Hilbert symbol computed in $\Q_p$, while $(a, b)_\infty$ is the Hilbert
symbol computed in $\R$.


\begin{theorem}
	\label{thm:almost_all_one}
	\lean{hilbertSym.almost_all_one}
	\uses{def:Hilbert_symbol}
	\leanok
Let $a, b \in \Q^\times$. Then $(a, b)_p = 1$ for all but finitely many primes $p$.
\end{theorem}

\begin{proof}
	TODO
\end{proof}

\begin{theorem}
	\label{thm:product_formula}
	\lean{hilbertSym.prod_eq_one}
	\uses{def:Hilbert_symbol}
	\leanok
Let $a, b \in \Q^\times$. Then $\prod_{v \in V} (a, b)_v = 1$.
\end{theorem}
\begin{proof}
	TODO
\end{proof}

\subsection{Approximation and Existence Theorem}

\begin{definition}
	\label{def:finite_embedding}
	\lean{Rat.finiteEmbedding}
	\leanok
Let $S$ be a finite set of prime numbers.
We define the \emph{finite embedding} of $\Q$ into $\R \times \prod_{p \in S} \Q_p$ as the diagonal map.

N.B. more precisely, in the Lean code, S is a subset of the natural numbers, and the product is
over the $n$-th prime for $n \in S$.
\end{definition}

\begin{theorem}\label{thm:approximation}%todo change references
	\lean{Rat.approximation}
	\uses{def:finite_embedding}
	\leanok
	Let $S$ be a finite set of prime numbers.
	The image of the finite embedding of $\Q$ into $\R \times \prod_{p \in S} \Q_p$ is dense with respect to
	the product topology. More concretely, for every $\epsilon > 0$ and every
	$y \in \R \times \prod_{p \in S} \Q_p)$, there exists $x \in \Q$ such that the distance
	between $y$ and the image of $x$ is less than $\epsilon$.
\end{theorem}
\begin{proof}
	TODO
\end{proof}


\begin{theorem}\label{thm:existence}
	\lean{hilbertSym.existence}
	\uses{def:Hilbert_symbol}
	\leanok
	Let $(a_i)_{i \in I}$ be a finite family of nonzero rational numbers, and let
	$e_{i,v} \in \{\pm 1\}$ for each $i \in I$ and each place $v$ of $\Q$. There exists a rational
	number $x \in \Q^\times$ such that
	\[
	\forall i \in I, \forall v \text{ place of } \Q, (x,a_i)_v=e_{i,v}
	\]
	if and only if the following conditions hold:
	\begin{itemize}
		\item For each $i \in I$, we have $e_{i,v}=1$ for all but finitely many places $v$.
		\item For each $i \in I$, we have $\prod_v e_{i,v}=1$.
		\item For each place $v$, there exists $x_v \in \Q_v^\times$ such that
		$(x_v, a_i)_v = e_{i,v}$ for all $i \in I$.
	\end{itemize}
\end{theorem}

The proof uses \code{Nat.infinite_setOf_prime_and_modEq}.

\begin{proof}
	\uses{def:Hilbert_symbol, thm:approximation}
TODO: necessity uses product formula, sufficiency uses the theorems in this section.
\end{proof}
